%% kapa: exam=2 round=21 year=2010 period=1 subject=appraisal_practice subject_ko=감정평가실무 seq=5 number=5 points=10
%% source: Q-Net 감정평가사 기출문제 (cst00309) artlSeq=5024183 file=extracted/1008/5024183/2010년 제21회 감정평가사 2차시험 기출문제/1교시.hwp
%% conversion: kordoc 4.24.51 -> tools/gen2.py (auto); review=proofread-sonnet; flags=none
\begin{longquestion}{【문제 5】}
주어진 자료는 2010년 상반기에 월별로 수집된 실거래 사례의 토지단가(원/m\textsuperscript{2})이다. A시 외곽의 동일수급권 내 자연녹지지역의 '답'에 대한 자료로서 용도 및 규모가 유사하며 제반 요인의 차이가 없다. 또한, 대상기간동안 지가변동도 미미하였다. 이 자료에만 의거하여 금년 7월 1일 기준으로 부동산가격공시 및 감정평가에 관한 법률상 언급되는 ‘성립될 가능성이 가장 높다고 인정되는 가격’을 결정하고자 한다. 다음의 순서에 입각하여 요구하는 값을 모두 구하고, 적정가격을 결정하되, 그 사유를 설명하시오. (10점) \par
물음(1) 범위(range) 및 평균(mean)의 산정 \par
물음(2) 중위값(median) 및 최빈치(mode)의 산정 \par
물음(3) 적정가격결정 사유의 설명 \par
<자료> 수집된 토지 가격자료 : 총 12개 \par
1) 2010년 1월 : 190,000원, 180,000원 \par
2) 2010년 2월 : 190,000원, 200,000원 \par
3) 2010년 3월 : 238,000원, 190,000원, 210,000원 \par
4) 2010년 4월 : 225,000원 \par
5) 2010년 5월 : 210,000원, 210,000원 \par
6) 2010년 6월 : 195,000원, 210,000원 \par
\end{longquestion}
